% Timing: 5:00--8:00   (paper: Sec. 1.1)
% One slide: the hard formula G_k, its DNF, and where the lower bound sits.
% (The DNF idea was already shown on the 'unavoidable?' slide; no section of its own.)

% ============================================================
\begin{frame}{The hard formula $G_k$}

\[
  G_k(x, y) \defeq \exists z \bigwedge_{i=1}^{k}
  \bigl[(P_i(x, y) \land Q_i(x, z)) \lor (Q_i(x, y) \land P_i(x, z))\bigr]
  \qquad \text{\NOTE{($|G_k| = O(k)$)}}
\]

\begin{center}
\begin{tikzpicture}[every node/.style={font=\small}]
  \node[pv,minimum size=7pt,label=left:$x$] (x) at (0,0) {};
  \node[pv,minimum size=7pt,label=right:$y$] (y) at (2.6,0) {};
  \node[pv,minimum size=7pt,label=above:$z$] (z) at (1.6,.6) {};
  \draw[thick,cA] (x) -- (y);
  \draw[thick,cB] (x) -- (z);
  \node[anchor=west,align=left] at (4.2,.6)
    {at every $i$: one of the edges $xy$, $xz$ is in $P_i$,\\
     the other is in $Q_i$};
\end{tikzpicture}
\end{center}

\vspace{1em}
DNF: \ $G_k \equiv \bigcup_{A \subseteq [k]} \bigl[ H_A \cap (H_{A^c} \mathbin{;} \one) \bigr]$
\ : \ \EMPH{$2^k$ compositions}
\ \NOTE{($H_A \defeq \bigcap_{i \in A} P_i \cap \bigcap_{i \notin A} Q_i$)}
\\ \hspace*{2.6em} $\Rightarrow$ optimal up to a factor of $2$: \ every equivalent \CoR{} term needs $\ge 2^{k-1}$ compositions


\note{
\textbf{スライド概要}：$G_k$ は DNF で合成 $2^k$ 個、これは 2 倍を除いて最適。

\textbf{メモ}：
\begin{itemize}
  \item 各 $i$ で、辺 $xy$ と辺 $xz$ の一方が $P_i$、他方が $Q_i$
  \item 下界は合成以外の演算（補・逆・共有）を無料にしても成り立つ
\end{itemize}
}
\end{frame}
